\documentclass[conference]{IEEEtran}
\IEEEoverridecommandlockouts
\usepackage{cite}
\usepackage{amsmath,amssymb,amsfonts}
\usepackage{algorithmic}
\usepackage{graphicx}
\usepackage{textcomp}
\usepackage{xcolor}
\usepackage{booktabs}
\usepackage{url}
\usepackage{hyperref}
\hypersetup{
    colorlinks=true,
    linkcolor=black,
    citecolor=black,
    urlcolor=blue,
}
\def\BibTeX{{\rm B\kern-.05em{\sc i\kern-.025em b}\kern-.08em
    T\kern-.1667em\lower.7ex\hbox{E}\kern-.125emX}}

\begin{document}

\title{Hierarchical Multi-Agent Control by LLM:\\
Simulation Environment, Per-Robot Navigation, and LLM Orchestration}

\author{
\IEEEauthorblockN{Dmitriy Vizitei}
\IEEEauthorblockA{\textit{Innopolis University} \\
Innopolis, Russia \\
d.vizitei@innopolis.university}
\and
\IEEEauthorblockN{Artyom Tuzov}
\IEEEauthorblockA{\textit{Innopolis University} \\
Innopolis, Russia \\
a.tuzov@innopolis.university}
\and
\IEEEauthorblockN{Dmitry Ryabov}
\IEEEauthorblockA{\textit{Innopolis University} \\
Innopolis, Russia \\
d.ryabov@innopolis.university}
}

\maketitle

\begin{abstract}
This report describes the Level-1 and Level-3 components of a
three-level hierarchical multi-agent system for natural-language
control of a 20-robot differential-drive swarm. Level~3 owns the
simulation environment and the per-robot navigation stack:
robot models, sensor pipelines, world files, transform tree,
and per-robot Nav2 lifecycle nodes. Level~1 is the LLM
orchestration layer that converts free-form natural-language
commands into concrete primitive invocations on the Level-2
swarm controller via three communication channels: a reactive
decision server, a passive observer, and an operator chat
interface. Our implementation is built on ROS~2 Humble inside
the shared \texttt{iros\_llm\_swarm} workspace. We deliberately
pivoted from Gazebo Harmonic to the 2D Stage simulator after
empirical evaluation showed that Gazebo could not sustain a
20-robot fleet at the per-robot complexity our project requires.
We document the simulation pivot and its rationale, describe the
per-robot Nav2 spawner and its parameter rewriting machinery,
present the five simulated scenarios including a dynamic obstacle
system and a richly-named map based on \emph{Among Us: The
Skeld}, and lay out the evaluation methodology planned for the
completed system.
\end{abstract}

\begin{IEEEkeywords}
multi-robot simulation, Stage, Gazebo, ROS~2, Nav2, LLM, tool
calling, behavior trees, dynamic obstacles
\end{IEEEkeywords}

\section{Introduction}

A hierarchical multi-agent system organizes control into a
small number of clearly-typed layers. In this project, three
levels are defined: Level~1, an LLM agent that interprets
free-form natural-language commands into structured tasks for
the fleet; Level~2, a classical swarm controller that turns
those tasks into conflict-free paths and stable formations;
and Level~3, the robots themselves, embedded in a simulated
environment that publishes ground-truth state, sensor data,
and consumes velocity commands. \emph{This report describes
the Level-1 and Level-3 contributions.} The Level-2 layer is
documented separately by the partner team.

The two ends of the hierarchy share a single concern that
makes them natural to discuss together: they both define the
\emph{interfaces} that the Level-2 layer is built against.
Level~3 fixes the topic graph the controller observes (every
\texttt{/robot\_N/odom}, every \texttt{/robot\_N/scan}, the
\texttt{/tf} tree, the static \texttt{/map}). Level~1 fixes
the action graph the controller is expected to expose
(\texttt{/swarm/set\_goals}, the \texttt{/formation/*}
services, the BT-node primitives). If either of those
contracts is wrong, the Level-2 layer cannot meet its
specification. The overall three-level structure is illustrated
in Fig.~\ref{fig:arch}.

\subsection*{Scope and contributions}

We focus on a 20-robot fleet of identical differential-drive
robots in five simulated scenarios. The Level-3 layer provides:
a deterministic, reproducible simulation environment that scales
to 20 simultaneously-spawned robots; plausible sensor and
odometry streams consumable by Nav2 and the upstream MAPF
planner; a per-robot Nav2 stack that comes up reliably under
DDS load; a dynamic obstacle system with runtime-togglable doors
and placeable obstacles; and a rich set of named map locations
that the Level-1 LLM can reference symbolically.

The Level-1 layer exposes a small, well-typed primitive surface
to the LLM; grounds free-form commands in the simulator's world
state without leaking low-level robot details; and integrates
with the existing behavior tree as the runtime substrate via
three distinct channels: a reactive decision server triggered
by BT/MAPF warnings, a passive observer that monitors fleet
state continuously, and an operator chat interface accessible
from RViz.

Our specific contributions are:
(1)~a documented pivot from Gazebo Harmonic to Stage 2D, with
the Gazebo configuration kept in tree for future use;
(2)~per-robot Nav2 spawning machinery using
\texttt{ReplaceString} and \texttt{RewrittenYaml} with
staggered launches to absorb DDS discovery cost;
(3)~a static \textsf{map} $\to$ \textsf{robot\_N/odom}
identity transform replacing AMCL at this fleet density;
(4)~five world files and associated map description YAMLs for
scenarios \textsf{cave}, \textsf{large\_cave},
\textsf{warehouse\_2}, \textsf{warehouse\_4}, and
\textsf{amongus} (\emph{Among Us: The Skeld});
(5)~a dynamic obstacle manager with runtime door and obstacle
APIs and scenario-level YAML configuration;
(6)~a three-channel LLM orchestration architecture
(decision server, passive observer, operator chat) grounded in
a behavior-tree interface over Level-2 primitives.

\section{Related Work}
\label{sec:related}

\subsection{Multi-robot simulation}

The two dominant simulation backends in the ROS ecosystem are
Stage \cite{vaughan2008stage} and Gazebo
\cite{koenig2004gazebo}, with Gazebo Harmonic
(\texttt{gz-sim}) being the modern split from the classic
Gazebo project. Stage is a lightweight 2D simulator built
explicitly for massively multi-robot scenarios, emphasizing
aggregate fleet throughput over per-robot fidelity. Gazebo
provides full 3D dynamics, rich sensor models, and physics
engines. The choice between them is a function of fleet size
and required fidelity: at small fleet sizes Gazebo's higher
fidelity is essentially free; at large fleet sizes Stage's
throughput is decisive. Stage's suitability for massive
multi-robot simulation is well documented
\cite{vaughan2008stage}. Recent generalised distributed
experimental grids such as Gengrid \cite{kedia2025gengrid}
target reproducible multi-robot experimentation.

\subsection{Differential-drive kinematics}

The differential-drive kinematic model is a well-studied
mobile-robot abstraction \cite{siwek2023diffdrive,siegwart2011}.
The robot is described by linear and angular velocity inputs
$(v, \omega)$ subject to nonholonomic constraints
$\dot{x} = v\cos\theta$, $\dot{y} = v\sin\theta$,
$\dot{\theta} = \omega$. Stage's differential drive model
implements this directly, with optional first-order lag on the
velocity inputs to mimic actuator dynamics. Recent work on
swarm-enabling technology \cite{chamanbaz2017swarmenabling}
demonstrates that this level of fidelity is sufficient for many
fleet-coordination questions when the upstream controller
already accounts for tracking error.

\subsection{ROS~2 and DDS at scale}

ROS~2 \cite{macenski2022ros2} replaces ROS~1's centralized
master with a peer-to-peer DDS middleware layer. At scale, DDS
discovery overhead becomes the dominant cost
\cite{maruyama2016dds,kronauer2024empirical}; tuning the RMW
implementation and pinning a profile is required to keep
multi-robot systems responsive. Nav2
\cite{macenski2020marathon,macenski2023nav2} is the de facto
ROS~2 navigation stack, with the layered costmap framework of
\cite{lu2014layered} as its obstacle-handling substrate.

\subsection{LLM-driven robotics}

A rapidly maturing line of work places large language models
in the loop with embodied agents.
SayCan~\cite{ahn2022saycan} grounds LLM proposals in a learned
affordance function. Code as Policies~\cite{liang2023codepolicies}
treats the LLM as a code generator over a fixed primitive API.
ProgPrompt~\cite{singh2023progprompt} extends this with situated
programs that include scene context. Inner
Monologue~\cite{huang2022inner} adds environmental feedback to
the prompt loop. RT-1~\cite{brohan2022rt1} and
RT-2~\cite{brohan2023rt2} take a different route, training
end-to-end models that output low-level control. For our
Level-1 layer the Code-as-Policies / SayCan / ProgPrompt
direction is the relevant one: we expose a fixed,
small-vocabulary primitive surface and ask the LLM to compose
it. Recent surveys of LLM agents in robotics
\cite{zeng2023llmrobsurvey} and tool-augmented language models
\cite{schick2023toolformer} cover the orchestration landscape.

\subsection{Behavior trees as an interface substrate}

Behavior trees \cite{colledanchise2018bt,iovino2022btsurvey}
compose atomic actions and conditions into a hierarchical
structure that is both executable and human-readable.
BehaviorTree.CPP \cite{btcpp} is a mature C++ implementation
that integrates with ROS~2 actions. For an LLM orchestration
layer, behavior trees offer two specific advantages: the tree
structure can be generated by the LLM, and runtime ticking is
handled by a deterministic engine rather than the LLM itself,
keeping latency and cost predictable.

\subsection{Multi-robot formation and fault tolerance}

The Level-2 contribution draws on
\cite{santos2020segregative,wang2025apf,jiang2020formation,martinez2023pso,debie2023swarmsurvey,elamvazhuthi2019meanfield,elsayed2024faulttolerant};
the partner team's report covers them in detail.

\section{Methodology}
\label{sec:method}

\subsection{Pivot from Gazebo to Stage}

The original proposal targeted Gazebo Harmonic (\texttt{gz-sim})
as the simulation backend. We implemented this configuration
end-to-end: URDF/Xacro robot descriptions, an SDF world file,
Nav2 bridges via \texttt{ros\_gz\_bridge}, and the same launch
orchestration we use today. The configuration is preserved in
the tree as \texttt{iros\_llm\_swarm\_simulation}.

Empirically, this configuration could not sustain a 20-robot
fleet on a single workstation. The dominant cost was the
simulator itself: physics step time, collision broadphase, and
sensor rendering scaled with fleet size in a way that pushed
the real-time factor well below usable thresholds before any
navigation stack was attached. With Nav2 also running, the
system was no longer viable for the coordination experiments
the project requires.

We therefore pivoted to Stage 2D \cite{vaughan2008stage} via
\texttt{stage\_ros2}. Stage's design assumption---many robots,
simple kinematics, simple sensors---is exactly aligned with the
project's needs. Our Stage configuration uses a unified TF tree
(\texttt{one\_tf\_tree: true}), enforced namespace prefixes
(\texttt{enforce\_prefixes: true}), a differential drive robot
model (\texttt{warehouse\_robot.inc}) with a planar laser
scanner, and five world files sharing the same robot model and
spawn protocol.

The Gazebo configuration is retained as the target for two
future activities: an aggressive performance pass (GPU-accelerated
sensors, simplified collision meshes, selective sensor-rate
degradation), and eventual hardware deployment where it serves as
a hardware-realistic preview environment.

\subsection{Simulation environment}

\paragraph{Robot model} The robot is a differential drive unit
defined by a footprint polygon, maximum linear and angular
velocities, and a planar laser sensor with configurable range
and angular resolution. Stage publishes
\texttt{nav\_msgs/Odometry} on \texttt{/robot\_N/odom} from
integration of commanded velocities (with optional Gaussian
noise), and \texttt{sensor\_msgs/LaserScan} on
\texttt{/robot\_N/scan} from raycasts against the world bitmap.

\paragraph{Scenarios and world files}
Five scenarios are defined in a shared
\texttt{common\_scenarios.yaml}. Each scenario specifies a
world file, a map YAML for Nav2, a map description YAML for the
LLM, 20 initial robot poses, and optional obstacle
configurations.

The \textsf{cave} scenario (75$\times$75~m) is a
narrow-passage dungeon-style map stressing single-file corridor
handling. \textsf{large\_cave} (150$\times$150~m) is the same
topology at double scale, intended for long-range path planning
experiments. \textsf{warehouse\_2} (30$\times$30~m) arranges
20 robots in two groups at opposite corners of a shelved
warehouse, while \textsf{warehouse\_4} places five robots in
each of the four corners for symmetric multi-group experiments.
The fifth scenario, \textsf{amongus}, maps the spaceship layout
of \emph{Among Us: The Skeld} (66.6$\times$37.75~m) with 14
named rooms (Cafeteria, Reactor, Navigation, Weapons, Storage,
MedBay, Admin, O$_2$, Security, Shields, Electrical,
Communications, Upper Engine, Lower Engine) and narrow corridor
choke-points that stress conflict resolution. The Among Us map
is included both as an engaging stress-test topology and as a
demonstration that the LLM can operate with fully symbolic
location references rather than raw coordinates.

\paragraph{Map descriptions for the LLM}
Each scenario ships a machine-readable map description YAML
consumed by the orchestration layer. These files enumerate
named locations with coordinates, location aliases in English
and Russian, colour-coded robot groups with home positions and
spawn coordinates, viable formation zones with free-space
radii, and planning heuristics. The heuristics encode
map-specific knowledge: bottleneck locations, unreachable goal
patterns, and guidance on when to abort versus replan. The LLM
is therefore able to reason about \emph{rooms}, \emph{groups},
and \emph{regions} rather than raw $(x,y)$ pairs.

\paragraph{TF tree}
A unified TF tree is enforced across all namespaces:
\[
\textsf{map}
  \to \textsf{robot\_N/odom}
  \to \textsf{robot\_N/base\_link}
\]
where the first link is a static identity transform published
by the per-robot Nav2 spawner (see
Section~\ref{sec:method:nav2}) and the second is published by
Stage. Sensor frames such as \textsf{robot\_N/laser\_link} are
children of \textsf{robot\_N/base\_link}.

\paragraph{Time source}
\texttt{use\_sim\_time=true} is the default everywhere. Stage
owns the simulated clock and publishes it on \texttt{/clock},
which is critical for timestamp-aware components such as
Nav2's costmaps and the MAPF schedule monitor.

\paragraph{Containerization}
The full development environment ships as a Docker image with
ROS~2 Humble, all dependencies, and the project workspace
pre-built. The image pins \texttt{ROS\_DOMAIN\_ID=42},
\texttt{RMW\_IMPLEMENTATION=rmw\_cyclonedds\_cpp}, and a
\texttt{CYCLONEDDS\_URI} pointing at the project DDS profile.
A \texttt{docker-compose.yaml} file provides standard service
definitions and X11 forwarding for RViz and Stage GUIs.

\subsection{Dynamic obstacle system}

The \texttt{iros\_llm\_swarm\_obstacles} package provides a
runtime-configurable obstacle manager. Two obstacle types are
supported: circles and axis-aligned rectangles (walls). In
addition, \emph{door} obstacles can be toggled open or closed
at runtime, allowing scenarios to model access-controlled
corridors. Obstacle layouts are loaded from scenario YAML at
launch time: \textsf{warehouse\_2} ships with a central
corridor door, while \textsf{amongus} ships with six
pre-placed corridor choke-point doors matching the map
geometry. New obstacles can also be placed interactively via a
custom RViz tool (\texttt{PlaceObstacleTool},
\texttt{DoorTool}) at any time during a run.

When an obstacle is added or a door state changes, the manager
recomputes the occupancy grid and publishes an updated
\texttt{/map}, triggering a MAPF replan on the Level-2 side.
This closes the loop between the LLM's ability to command
obstacle state changes and the planner's ability to respond.

\subsection{Per-robot Nav2 stack}
\label{sec:method:nav2}

Each robot runs its own Nav2 stack:
\texttt{controller\_server} (DWB local trajectory tracking),
\texttt{behavior\_server} (recovery actions), a
\texttt{lifecycle\_manager} that brings the above to the
\textsf{active} state, and a static
\textsf{map}$\to$\textsf{robot\_N/odom} identity transform
publisher. A single shared \texttt{map\_server} serves all
robots; the static map is namespace-independent.

\paragraph{Parameter rewriting}
The Nav2 parameter file is written once, with per-namespace
fields parameterised by template strings. Two facilities are
used to expand them at launch time. \texttt{ReplaceString}
substitutes the robot namespace into hard-coded occurrences of
the placeholder. \texttt{RewrittenYaml} re-keys the YAML file
at the top level and overrides specific values such as
\texttt{use\_sim\_time}. The combination produces, from a
single template, a fully-namespaced parameter file for every
robot.

\paragraph{Costmap configuration}
The local costmap is configured with the project-specific
\texttt{ResettingObstacleLayer} (Level-2 contribution) and an
inflation layer. The \texttt{observation\_persistence}
parameter is set to \texttt{0.0} so the costmap does not
retain stale lethal zones once a robot has moved; the local
costmap radius is chosen to exceed
\texttt{obstacle\_max\_range}.

\paragraph{Static localization}
AMCL is intentionally disabled. With 20 robots in a shared
workspace, every robot's laser scan contains the bodies of
every other robot, which are not reflected in the static map;
the particle filter degrades. We replace AMCL with a static
identity transform from \textsf{map} to
\textsf{robot\_N/odom}, which is correct in simulation since
the simulator's odometry is ground-truth-aligned. On hardware,
this identity transform must be replaced by a multi-robot
localization or SLAM pipeline.

\paragraph{Staggered bring-up}
Per-robot Nav2 instances are launched with a $\sim$0.3~s
stagger to absorb DDS discovery cost. Without stagger,
simultaneous launches reliably miss discovery on a fraction of
the robots and the system comes up incomplete.

\subsection{LLM orchestration layer}

The Level-1 LLM agent is implemented as the
\texttt{iros\_llm\_orchestrator} package. It communicates with
the rest of the system via three channels, each serving a
different interaction pattern. The underlying model is
configurable at launch time via a shared YAML; the deployed
configuration targets Llama-3.3-70B served through a
Groq-hosted OpenAI-compatible endpoint, with temperature 0.1
to suppress non-determinism.

\paragraph{Channel 1 --- reactive decision server}
The \texttt{llm\_decision\_server} node receives
\texttt{/llm/decision} requests from the behavior tree
whenever a BT or MAPF node emits a \texttt{WARN} or
\texttt{ERROR} status. Given the current fleet state---mode,
active formation, recent BT events, and the map
description---the server returns one of three decisions:
\texttt{wait}, \texttt{abort}, or \texttt{replan}. The
decision is logged to a per-run dataset at
\texttt{\~{}/.ros/llm\_decisions} for offline analysis. At
most one decision is processed concurrently; concurrent
requests queue behind it.

\paragraph{Channel 2 --- passive observer}
The \texttt{llm\_passive\_observer} node subscribes to
\texttt{/bt/state} and monitors fleet status continuously.
When a \texttt{WARN} or \texttt{ERROR} condition persists
beyond a configurable cooldown, the observer autonomously
composes a recovery command and publishes it on
\texttt{/llm/command}, bypassing the BT's decision request
flow. This channel is disabled by default and activated with
the \texttt{enable\_passive\_observer:=true} launch argument.
A dataset of autonomous interventions is logged to
\texttt{\~{}/.ros/llm\_commands}.

\paragraph{Channel 3 --- operator chat}
The \texttt{llm\_chat\_server} node exposes a conversational
interface on \texttt{/llm/chat}. The operator types a
natural-language command in the RViz panel; the server
assembles a context packet (BT state, active formations, map
summary, recent event history) and sends it together with the
command to the LLM, which returns a structured JSON plan that
is dispatched as a \texttt{/swarm/set\_goals} action or a
formation service call. Context is optionally populated via an
MCP bridge (\texttt{ros-mcp} over stdio), which provides
read-only access to live ROS topics without requiring the
orchestrator to subscribe to them directly.

\paragraph{Primitive surface}
All three channels operate over the same underlying primitive
surface exposed by Level-2: the \texttt{/swarm/set\_goals}
action, the \texttt{/formation/*} CRUD services, and the
\texttt{/obstacles/*} API introduced by this team.
Corresponding behavior-tree action nodes are
\texttt{MapfPlan}, \texttt{SetFormation},
\texttt{DisableFormation}, \texttt{CheckMode}, and
\texttt{ObstaclesMode}. The LLM is not permitted to publish
raw \texttt{cmd\_vel} or per-robot paths; the abstraction
boundary is enforced by the BT engine.

\paragraph{BT as the substrate}
The behavior tree is the runtime substrate over which the LLM
operates. The LLM does not directly call ROS~2 actions;
instead, it generates a behavior-tree XML (or selects from a
template library) whose action nodes are the primitive
wrappers listed above. The BT is then ticked by the
deterministic \texttt{BehaviorTree.CPP}~v4 engine. This
decomposition has three consequences: the LLM's output space
is structured XML against a known schema and is easy to
validate; the LLM is invoked only at decision boundaries, not
on every tick; and behavior-tree fallbacks and guards encode
safety properties that the LLM cannot bypass even if its
output is degenerate.

\paragraph{Grounding and state context}
The LLM receives a compact summary of fleet state derived from
\texttt{FollowerStatus}, \texttt{FormationsStatus}, and the
aggregated BT event stream. Raw scans and per-cell costmap
data are deliberately excluded: the LLM is expected to reason
about robots, formations, and named regions of the map, not
laser pixels. The map description YAML (Section~\ref{sec:method})
is included verbatim in the system prompt, giving the LLM
access to named locations, viable formation zones, and
planning heuristics without any runtime lookup.

\section{Code Availability}
\label{sec:github}

The full implementation, shared with the Level-2 team, is
hosted at:
\begin{center}
\url{https://github.com/innopolis-robotics-society/llm_swarm}
\end{center}
The repository contains a top-level \texttt{README.md} with
Docker bring-up instructions, scenario launch commands, and a
description of the goal JSON format. Per-package documentation
is generated by \texttt{rosdoc2} and aggregated into a Sphinx
site under \texttt{iros\_llm\_swarm\_docs/}.

The packages owned by this team are
\texttt{iros\_llm\_swarm\_simulation\_lite} (Stage~2D worlds,
robot model, scenario definitions, map descriptions),
\texttt{iros\_llm\_swarm\_simulation} (Gazebo Harmonic
preview), \texttt{iros\_llm\_swarm\_local\_nav} (per-robot
Nav2 spawner), \texttt{iros\_llm\_swarm\_obstacles} (dynamic
obstacle manager and RViz tools), \texttt{iros\_llm\_orchestrator}
(three-channel LLM orchestration), and
\texttt{iros\_llm\_rviz\_panel} (operator chat panel).

\section{Evaluation}
\label{sec:exp}

All tests run inside the project's Docker image on a single
workstation with \texttt{ROS\_DOMAIN\_ID=42},
\texttt{RMW\_IMPLEMENTATION=rmw\_cyclonedds\_cpp}, and the
project CycloneDDS profile.

\subsection{Unit and integration tests}

The orchestration layer is covered by a suite of offline
Python tests that run without a live ROS graph.
\texttt{test\_decision\_parser} verifies that
\texttt{parse\_llm\_decision()} correctly extracts the
\texttt{wait}/\texttt{abort}/\texttt{replan} verdict from
plain JSON, fenced JSON, prose-wrapped JSON, and degenerate
inputs (empty string, wrong type, unknown enum value); all
cases fall back to \texttt{wait} on malformed input.
\texttt{test\_decision\_server} spins a
\texttt{LlmDecisionServer} in mock mode in a background thread
and fires action-client requests over loopback, verifying that
the correct decision is returned, feedback stages
(\texttt{received}/\texttt{thinking}/\texttt{done}) are
observed in order, and a JSONL dataset entry is written per
call. \texttt{test\_passive\_observer} stands up a fake
\texttt{/llm/command} action server and verifies that a
sustained \texttt{WARN} \texttt{BTState} message triggers
exactly one proactive goal, while a second \texttt{WARN}
within the cooldown window is suppressed.
\texttt{test\_prompt\_builder} and
\texttt{test\_command\_parser} cover the structured-output
parsing path end-to-end with fixtures that exercise all
supported plan fields and edge cases.

The MAPF solver is validated by two compiled stress tests.
A smoke test places three agents on crossing straight paths
in a 20$\times$20 empty grid and checks that
\texttt{LNS2Solver::solve()} returns \texttt{ok=true} with
zero final collisions. A corridor stress test places eight
agents in opposing pairs that must share a four-cell-wide
corridor in a 30$\times$30 grid with blocked interior; the
solver must resolve all head-on conflicts within a 3~s
budget. Both tests pass deterministically at the fixed seeds
used in CI.

The behavior tree is exercised by \texttt{test\_bt\_runner},
a standalone ROS~2 node that runs the full
\texttt{FleetMain} tree against a live 20-robot Stage
simulation without an LLM in the loop. A scripted six-step
scenario is executed: (1)~all 20 robots converge to the
warehouse centre in a $4{\times}5$ grid; (2)~a wedge
formation is activated over the orange squad; (3)~a line
formation is activated over the blue squad; (4)~a full
cross-swap sends orange to the top-right and blue to the
bottom-left, stressing MAPF conflict resolution over the
full 30~m diagonal; (5)~all robots return home; (6)~idle.
The scenario completes successfully in the warehouse
scenario with staggered DDS bring-up.

\subsection{Scenario bring-up}

All five world presets (\textsf{cave}, \textsf{large\_cave},
\textsf{warehouse\_2}, \textsf{warehouse\_4},
\textsf{amongus}) bring up cleanly with the full stack via
\texttt{swarm\_lns.launch.py}. All 20 robots reach the
\textsf{active} Nav2 lifecycle state within the staggered
bring-up window. The MAPF planner produces valid plans in
each scenario; the \textsf{large\_cave} scenario requires
longer initial planning time due to map scale but does not
time out. In the \textsf{amongus} scenario, named-location
references (\emph{Cafeteria}, \emph{Reactor}, etc.) resolve
correctly via the map description YAML, and toggling a
corridor door triggers a MAPF replan within one planning
cycle.

\subsection{Limitations}

Several aspects of the system have not been subjected to
systematic quantitative evaluation and constitute open items
for future work.

\textit{Simulator throughput.} The Gazebo-vs-Stage throughput
comparison---real-time factor as a function of $N$ and
navigation load---has not been measured beyond the qualitative
observation that motivated the pivot. A controlled sweep
across $N \in \{1,5,10,15,20\}$ with and without Nav2
attached would quantify the crossover point and the headroom
available per robot in the Stage configuration.

\textit{Nav2 bring-up reliability at scale.} We observe that
staggered launch reliably brings up all 20 robots, but we
have not collected statistics over repeated cold boots.
Failure rates as a function of stagger interval and DDS
profile have not been measured.

\textit{LLM decision quality.} The decision server and
passive observer have been tested in mock mode and with
live LLM calls in ad-hoc sessions, but no held-out incident
dataset with ground-truth labels has been assembled. Metrics
such as decision accuracy, schema validity rate, and
command-to-execution latency remain unmeasured.

\textit{End-to-end command grounding.} The full pipeline from
natural-language operator command to completed fleet mission
has been demonstrated interactively but not evaluated on a
systematic command set. Coverage across the five scenarios
and the three orchestration channels is therefore anecdotal.

\section{Analysis and Observations}
\label{sec:analysis}

\paragraph{Gazebo did not scale to our fleet}
The pivot to Stage was driven by the simulator alone, not by
the algorithms running on top. Gazebo's per-robot
cost---collision broadphase, sensor rendering, physics solver
work---scales steeply with fleet size, and a 20-robot
configuration on a single workstation leaves no usable
headroom for navigation, planning, and behavior-tree ticking.
The corollary is that the Gazebo configuration is not wrong;
it is correct for a smaller fleet or a multi-machine
deployment. The pivot is recorded as such, not as a deletion.

\paragraph{Static localization is required, not preferred}
We attempted several AMCL configurations before disabling it.
The fundamental issue is that 19 of the 20 laser scans visible
to any given robot contain the bodies of other robots, which
the static map does not predict; the particle filter's
likelihood field collapses onto incorrect hypotheses. The
static identity transform replaces this with a
known-correct fact (the simulator's odometry is
ground-truth-aligned). The system will require a multi-robot
SLAM solution before any hardware deployment.

\paragraph{Parameter rewriting is a fragile point}
Per-namespace Nav2 parameter files are produced from templates
by \texttt{ReplaceString} and \texttt{RewrittenYaml}. A
specific failure mode required several days to diagnose: the
\texttt{use\_sim\_time} key, when nested incorrectly under the
\texttt{root\_key} introduced by \texttt{RewrittenYaml},
silently fails to propagate to the costmap layer. The costmap
then runs on wall-clock timestamps while the rest of the stack
uses simulated time; every laser scan is rejected as too old,
the local planner produces no obstacles, and the costmap
remains empty with no explicit error. Any Nav2 deployment
using \texttt{RewrittenYaml} should validate
\texttt{use\_sim\_time} propagation as a first-class test.

\paragraph{Costmap radius constraints are silent}
A second silent failure: if the local costmap's physical radius
is smaller than \texttt{obstacle\_max\_range}, laser returns
at long range fall outside the costmap window and are dropped
at the range check, leaving a costmap that contains only
short-range obstacles. The robot drives past mid-range
obstacles because, from its costmap's perspective, they do not
exist.

\paragraph{Minimal Nav2 needs both planner and BT}
An early configuration ran only \texttt{controller\_server},
on the hypothesis that the Level-2 planner would supply
complete paths and the local controller would track them. This
configuration is unstable: DWB relies on a global path from
\texttt{planner\_server}, and recovery behaviors are only
available through \texttt{bt\_navigator}. Nav2 expects both,
even when the global path is externally authoritative.

\paragraph{DDS is the dominant bring-up cost}
On the wall-clock budget of a 20-robot launch, DDS discovery
is larger than physics step time, larger than Nav2
initialization, and larger than the planner's first plan.
Without the project DDS profile, a non-trivial fraction of
bring-up runs fails to register all robots within the
staggered timeline.

\paragraph{The LLM interface is not a free choice}
Choosing to expose the system to the LLM at the level of the
behavior tree---rather than at the level of individual actions
or raw \texttt{cmd\_vel}---reflects a specific architectural
bet: the LLM should be a slow, infrequent reasoner over
semantic primitives, not a tight controller. This matches both
the latency profile of current frontier models and the safety
properties required for 20-robot control. The three-channel
structure (reactive, passive, and operator-driven) allows the
system to mix LLM-initiated and human-initiated commands
without a unified arbitration mechanism at the BT level.

\paragraph{Named map locations improve command quality}
Early testing with raw coordinates in LLM prompts produced a
higher rate of out-of-bounds and wall-interior goal
coordinates than symbolic references. Providing the
\textsf{named\_locations} and \textsf{formation\_zones}
blocks in the system prompt, together with planning heuristics
specific to each map, substantially reduced goal validation
failures on the Level-2 side. The \textsf{amongus} map, with
its 14 named rooms and rich alias table, is the clearest
demonstration of this principle.

\section{Conclusion and Future Work}
\label{sec:conclusion}

We described the Level-1 and Level-3 layers of a hierarchical
multi-robot control stack: a Stage 2D simulation environment
chosen after an explicit and documented pivot from Gazebo
Harmonic; five simulated scenarios including a dynamic
obstacle system and a richly-named map; a per-robot Nav2 stack
with parameter rewriting and staggered bring-up suitable for
20-robot fleets; a static localization strategy replacing AMCL
at this density; and a three-channel LLM orchestration
architecture grounded in a behavior-tree interface over
Level-2 primitives.

Future work falls into three directions. First,
\emph{simulator return path}: the preserved Gazebo Harmonic
configuration is the target for a performance pass (GPU
sensors, simplified collision geometry, selective sensor
degradation) that would restore 3D simulation at this fleet
size, and ultimately for hardware deployment where physical
robots replace the simulator entirely. Second,
\emph{multi-robot localization}: the static identity
transform from \textsf{map} to \textsf{robot\_N/odom} is
correct in simulation but must be replaced by a multi-robot
SLAM or inter-robot range-sharing pipeline before any hardware
transfer. Third, \emph{systematic evaluation of the LLM
layer}: assembling held-out incident and command datasets to
measure decision accuracy, schema validity rate, and
end-to-end grounding quality across all three orchestration
channels and all five scenarios, as described in
Section~\ref{sec:exp}.

\begin{thebibliography}{00}

\bibitem{vaughan2008stage}
R. Vaughan,
``Massively multi-robot simulation in Stage,''
\emph{Swarm Intell.}, vol. 2, pp. 189--208, 2008.

\bibitem{koenig2004gazebo}
N. Koenig and A. Howard,
``Design and use paradigms for Gazebo, an open-source
multi-robot simulator,''
in \emph{Proc. IROS}, 2004.

\bibitem{kedia2025gengrid}
P. Kedia and M. Rao,
``Gengrid: A generalised distributed experimental
environmental grid for swarm robotics,''
arXiv:2504.20071, 2025.

\bibitem{siwek2023diffdrive}
M. Siwek, J. Panasiuk, L. Baranowski, W. Kaczmarek,
P. Prusaczyk, and S. Borys,
``Identification of differential drive robot dynamic model
parameters,''
\emph{Materials}, vol. 16, no. 2, p. 683, 2023.

\bibitem{siegwart2011}
R. Siegwart, I. R. Nourbakhsh, and D. Scaramuzza,
\emph{Introduction to Autonomous Mobile Robots}, 2nd ed.
Cambridge, MA, USA: MIT Press, 2011.

\bibitem{chamanbaz2017swarmenabling}
M. Chamanbaz \emph{et al.},
``Swarm-enabling technology for multi-robot systems,''
\emph{Front. Robot. AI}, vol. 4, p. 12, 2017.

\bibitem{macenski2022ros2}
S. Macenski, T. Foote, B. Gerkey, C. Lalancette, and
W. Woodall,
``Robot Operating System 2: Design, architecture, and uses
in the wild,''
\emph{Sci. Robot.}, vol. 7, no. 66, 2022.

\bibitem{maruyama2016dds}
Y. Maruyama, S. Kato, and T. Azumi,
``Exploring the performance of ROS2,''
in \emph{Proc. EMSOFT}, 2016.

\bibitem{kronauer2024empirical}
T. Kronauer, J. Pohlmann, M. Matthe, T. Smejkal, and
G. Fettweis,
``Latency analysis of ROS2 multi-node systems,'' 2023.

\bibitem{macenski2020marathon}
S. Macenski, F. Mart\'in, R. White, and J. Gin\'es Clavero,
``The Marathon 2: A navigation system,''
in \emph{Proc. IROS}, 2020.

\bibitem{macenski2023nav2}
S. Macenski \emph{et al.},
``From the desks of ROS maintainers: A survey of modern \&
capable mobile robotics algorithms in the robot operating
system 2,''
\emph{Robot. Auton. Syst.}, 2023.

\bibitem{lu2014layered}
D. V. Lu, D. Hershberger, and W. D. Smart,
``Layered costmaps for context-sensitive navigation,''
in \emph{Proc. IROS}, 2014.

\bibitem{ahn2022saycan}
M. Ahn \emph{et al.},
``Do as I can, not as I say: Grounding language in robotic
affordances,''
in \emph{Proc. CoRL}, 2022.

\bibitem{liang2023codepolicies}
J. Liang \emph{et al.},
``Code as policies: Language model programs for embodied
control,''
in \emph{Proc. ICRA}, 2023.

\bibitem{singh2023progprompt}
I. Singh \emph{et al.},
``ProgPrompt: Generating situated robot task plans using
large language models,''
in \emph{Proc. ICRA}, 2023.

\bibitem{huang2022inner}
W. Huang \emph{et al.},
``Inner monologue: Embodied reasoning through planning with
language models,''
in \emph{Proc. CoRL}, 2022.

\bibitem{brohan2022rt1}
A. Brohan \emph{et al.},
``RT-1: Robotics transformer for real-world control at
scale,'' arXiv:2212.06817, 2022.

\bibitem{brohan2023rt2}
A. Brohan \emph{et al.},
``RT-2: Vision-language-action models transfer web knowledge
to robotic control,'' arXiv:2307.15818, 2023.

\bibitem{schick2023toolformer}
T. Schick \emph{et al.},
``Toolformer: Language models can teach themselves to use
tools,''
in \emph{Proc. NeurIPS}, 2023.

\bibitem{zeng2023llmrobsurvey}
F. Zeng \emph{et al.},
``Large language models for robotics: A survey,'' 2023.

\bibitem{colledanchise2018bt}
M. Colledanchise and P. {\"O}gren,
\emph{Behavior Trees in Robotics and AI: An Introduction}.
CRC Press, 2018.

\bibitem{iovino2022btsurvey}
M. Iovino, E. Scukins, J. Styrud, P. {\"O}gren, and C. Smith,
``A survey of behavior trees in robotics and AI,''
\emph{Robot. Auton. Syst.}, vol. 154, 2022.

\bibitem{btcpp}
D. Faconti and M. Colledanchise,
``BehaviorTree.CPP,''
\url{https://www.behaviortree.dev}, accessed 2026.

\bibitem{santos2020segregative}
V. G. Santos \emph{et al.},
``Spatial segregative behaviors in robotic swarms using
differential potentials,''
\emph{Swarm Intelligence}, vol. 14, no. 4, pp. 259--284, 2020.

\bibitem{wang2025apf}
H. Wang, M. Yue, X. Sun, and X. Zhao,
``Cooperative formation control strategy for multi-robot
system based on APF algorithm and sliding mode estimator,''
\emph{Robot. Auton. Syst.}, vol. 193, p. 105075, 2025.

\bibitem{jiang2020formation}
C. Jiang, Z. Chen, and Y. Guo,
``Multi-robot formation control: a comparison between
model-based and learning-based methods,''
\emph{J. Control Decis.}, vol. 7, no. 1, pp. 90--108, 2020.

\bibitem{martinez2023pso}
F. Martinez and A. Rendon,
``Autonomous motion planning for a differential robot using
particle swarm optimization,''
\emph{Int. J. Adv. Comput. Sci. Appl.}, vol. 14, 2023.

\bibitem{debie2023swarmsurvey}
E. Debie, K. Kasmarik, and M. Garratt,
``Swarm robotics: A survey from a multi-tasking
perspective,''
\emph{ACM Comput. Surv.}, vol. 56, no. 2, 2023.

\bibitem{elamvazhuthi2019meanfield}
K. Elamvazhuthi and S. Berman,
``Mean-field models in swarm robotics: A survey,'' 2019.

\bibitem{elsayed2024faulttolerant}
A. M. Elsayed, M. Elshalakani, S. A. Hammad, and S. A. Maged,
``Decentralized fault-tolerant control of multi-mobile robot
system addressing LiDAR sensor faults,''
\emph{Sci. Rep.}, vol. 14, no. 1, p. 25713, 2024.

\bibitem{li2021lns2}
J. Li, Z. Chen, D. Harabor, P. J. Stuckey, and S. Koenig,
``Anytime multi-agent path finding via large neighborhood
search,''
in \emph{Proc. IJCAI}, 2021.

\end{thebibliography}

\end{document}